\documentclass[11pt]{article}

% Basic packages
\usepackage[utf8]{inputenc}
\usepackage{amsmath,amssymb,amsthm,amsfonts}
\usepackage{hyperref}
\usepackage{geometry}
\usepackage{enumitem}
\usepackage{bm}
\usepackage{mathtools}
\usepackage{graphicx}
\usepackage{color}
\usepackage{titlesec}
\usepackage{cite}
\usepackage{mathrsfs}
\usepackage{courier}
\usepackage{listings}

\geometry{margin=1in}
\setlength{\parskip}{0.6em}
\setlength{\parindent}{0em}

\titleformat{\section}{\large\bfseries}{\thesection}{0.6em}{}
\titleformat{\subsection}{\normalsize\bfseries}{\thesubsection}{0.6em}{}
\titleformat{\subsubsection}{\normalsize\itshape}{\thesubsubsection}{0.6em}{}

\title{Non-Associative Fractal Learning for Quantum AGI:\\
Prime-Indexed Tensor Networks and Path-Dependent Intelligence}

\author{Citizen Gardens \\ The Foundation of Multiplicity}
\date{April 25, 2026}

% Simple math macros
\newcommand{\R}{\mathbb{R}}
\newcommand{\N}{\mathbb{N}}
\newcommand{\Mcal}{\mathcal{M}}
\newcommand{\Acal}{\mathcal{A}}
\newcommand{\Fcal}{\mathcal{F}}
\newcommand{\Lcal}{\mathcal{L}}
\newcommand{\E}{\mathbb{E}}
\newcommand{\Prob}{\mathbb{P}}
\newcommand{\norm}[1]{\left\lVert #1 \right\rVert}

% Multiplicity / QAGI-specific macros
\newcommand{\Lm}{\Lambda_m}
\newcommand{\Ap}{A_p}
\newcommand{\Aw}{\Acal_{w}}
\newcommand{\Awt}{\Acal_{w(t)}}
\newcommand{\Awp}{\Acal_{w(t+1)}}

% Listings setup
\lstdefinestyle{code}{
    basicstyle=\ttfamily\small,
    breaklines=true,
    frame=single,
    columns=flexible,
    captionpos=b
}
\lstset{style=code}

\begin{document}
\maketitle

\begin{abstract}
This document presents a comprehensive description of \emph{Non-Associative Fractal Learning} (NAFL), a cognitive and mathematical architecture for Quantum Agentic Governance Intelligence (QAGI) based on Prime-Indexed Tensor Networks (PITNs) and Multiplicity Theory.
Unlike classical deep learning, which compresses history into flat latent vectors, NAFL uses non-associative prime-indexed operators to make the sequence of operations---the \emph{prime word}---the primary carrier of memory and structure.
We formalize the operator algebra, the induced fractal topology of the knowledge manifold, the role of the Universal Multiplicity Constant \(\Lm\) as a stabilizer, and the construction of hierarchical decision policies \(\pi\) that operate over composite operators rather than static states.
We also contrast NAFL with classical reinforcement learning and hierarchical RL, and provide a minimal code-style specification of a toy implementation to support future empirical work and to establish technical prior art.
\end{abstract}
\newpage
\tableofcontents

\section{Executive Summary}

Non-Associative Fractal Learning (NAFL) is a learning and decision-making framework designed for advanced artificial general intelligence within the QAGI program.
The central idea is that memory and intelligence are not properties of a flat state vector, but of an evolving \emph{path} in a non-associative operator algebra indexed by prime numbers.

In NAFL, each primitive cognitive or computational act is implemented as a prime-indexed operator \(A_p\) acting on a tensorial manifold \(\Mcal\).
Instead of imposing associativity, the composition of these operators is deliberately non-associative, so that
\[
(A_{p_2} A_{p_3}) A_{p_5} \neq A_{p_2} (A_{p_3} A_{p_5}),
\]
and the ordered, bracketed sequence of operators---the \emph{prime word}---is itself the semantic identity of the resulting transformation.
This makes the system intrinsically path-dependent and encodes ``how'' a conclusion was reached as part of the state, rather than discarding it.

The induced state space has a fractal, tree-like topology: each new prime application extends a branch, and the same finite alphabet of primes is reused across scales.
A meta-learning parameter \(\beta(t)\) controls which primes are available at a given stage of development, leading to dimensional unfolding as the system encounters novel environments.
The Universal Multiplicity Constant \(\Lm\) acts as a stabilizer, ensuring that recursive operator application remains bounded and lawfully convergent.

Decision-making is formulated as a policy \(\pi\) that chooses the next prime operator to extend a word \(w(t)\) so that the resulting composite operator \(\Awp\) lands in high-reward, stable regions of \(\Mcal\).
This supports hierarchical, contextual decision-making, resilience to catastrophic forgetting, and cross-domain ecological awareness via an entanglement functional \(Q_{\mathrm{ent}}\) that couples structurally related branches of the operator tree.

This report consolidates the conceptual and mathematical foundations of NAFL, offers precise definitions, and sketches a minimal implementation pathway.
It is intended both as a working reference for continued development and as prior art establishing the core ideas: prime-indexed non-associative operator algebras, fractal knowledge manifolds, and path-dependent decision policies for AGI.

\section{Background and Motivation}

\subsection{Limitations of classical deep learning and RL}

Classical deep learning architectures typically represent information as vectors in high-dimensional Euclidean or Riemannian manifolds, with learning implemented as a sequence of associative linear and nonlinear transforms.
Although recurrent architectures and attention mechanisms can preserve some sequential information, the learned weights themselves usually implement effectively associative composition at inference time: the ordering of independent layers is fixed and does not encode a dynamically evolving algebra.

Reinforcement learning (RL) methods, particularly in their standard Markov decision process (MDP) setting, assume that the current state \(s_t\) contains all relevant information.
History is either discarded or compressed into a hidden state that is not structurally represented as a compositional object.
Hierarchical RL (HRL) introduces options and temporally extended actions, but these hierarchies are often manually designed or discovered through heuristic methods.

This combination of flat state representation and associative computation makes it difficult to:
\begin{itemize}[nosep]
    \item represent detailed path dependence,
    \item build emergent, scale-consistent hierarchies without hand-crafting, and
    \item avoid catastrophic forgetting when new tasks are introduced.
\end{itemize}

\subsection{Multiplicity Theory and QAGI}

Multiplicity Theory proposes that the recurrence and decomposition of elements through prime structures is foundational for modeling nonlinear, emergent phenomena across mathematics, physics, computation, and cognition.
In this view, mathematical objects are not static entities but recursively generated patterns of prime-labeled interactions.
Sets and modules become multiplicity spaces with identity and behavior preserved through recursive feedback across scales.

The QAGI (Quantum AGI) framework builds on this perspective by introducing Prime-Indexed Tensor Networks (PITNs) as physical and computational substrates.
Each prime index labels a distinct ``channel'' or mode of interaction, and the overall dynamics are governed by a universal constant \(\Lm\), which stabilizes evolution and encodes lawful constraints across scales.

NAFL emerges as the cognitive architecture that sits on top of PITNs, specifying how prime-indexed operators are composed, how the resulting state space is navigated, and how decisions are made.

\section{Prime-Indexed Tensor Networks and Non-Associativity}

\subsection{Prime-indexed operators}

Let \(\mathbb{P} = \{2,3,5,7,\dots\}\) denote the set of prime numbers.
To each prime \(p \in \mathbb{P}\) we associate an operator
\[
A_p : \Mcal \to \Mcal,
\]
where \(\Mcal\) is a tensor manifold representing the internal configuration space of the QAGI system.
Concretely, \(\Mcal\) may be realized as a manifold of tensor network states (e.g. MPS, TTN, PEPS) of fixed or variable bond dimension.

The family \(\{A_p\}_{p \in \mathbb{P}}\) forms a generating set for a non-associative algebra of transformations.
The operators may be parameterized (learned) maps, subject to architectural constraints (e.g., implemented by tensor network layers) and global stability constraints controlled by \(\Lm\).

\subsection{Non-associative composition}

Given two primes \(p,q \in \mathbb{P}\), the composition \(A_q A_p\) denotes the operator that applies \(A_p\) first, then \(A_q\):
\[
A_q A_p(x) := A_q\big( A_p(x) \big), \quad x \in \Mcal.
\]
We explicitly do \emph{not} impose associativity; instead, for three primes \(p,q,r\) we have in general
\[
(A_r A_q) A_p \neq A_r (A_q A_p).
\]
The space of all finite compositions of \(\{A_p\}\), with the bracketing structure preserved, forms a free (or constrained) non-associative algebra.
Each composite operator is uniquely identified by both:
\begin{itemize}[nosep]
    \item the ordered sequence of primes \((p_{i_1}, \dots, p_{i_t})\), and
    \item the specific bracketing (parenthesization) of the composition.
\end{itemize}

In the present work we focus on a fixed bracketing convention, typically right-associative:
\[
\Acal_{w(t)} := A_{p_{i_t}} A_{p_{i_{t-1}}} \dots A_{p_{i_1}},
\]
with the understanding that non-associativity manifests in deeper variants where the bracketing itself is a degree of freedom.

\subsection{Prime words as memory}

Define a \emph{prime word} of length \(t\) as
\[
w(t) = (p_{i_1}, p_{i_2}, \dots, p_{i_t}), \quad p_{i_k} \in \mathbb{P}.
\]
To each word \(w(t)\) we associate the composite operator
\[
\Acal_{w(t)} = A_{p_{i_t}} A_{p_{i_{t-1}}} \dots A_{p_{i_1}}.
\]
Given an initial configuration \(x_0 \in \Mcal\), the state at time \(t\) is
\[
x_t = \Acal_{w(t)}(x_0).
\]

In NAFL, the pair \((x_t, w(t))\) is the true state.
Crucially, two states with identical \(x_t\) but different prime words are treated as distinct because the future evolutions they admit are different: the available operators, their local spectra, and their entanglement relations depend on the entire path.

Thus, memory is not an auxiliary buffer but is \emph{intrinsic} to the algebra of transformations.
The path \(w(t)\) is a first-class citizen in the state space.

\subsection{Stabilization via the Universal Multiplicity Constant}

The Universal Multiplicity Constant \(\Lm\) is a scalar (or small collection of scalars) that encodes global stability and lawfulness of recursive dynamics.
We model its role by defining scaled operators
\[
\widetilde{A}_p := \Lm A_p,
\]
with \(\Lm\) chosen so that composites remain bounded in an appropriate operator norm:
\[
\norm{\widetilde{\Acal}_{w(t)}} \leq C(\Lm)
\]
for all words \(w(t)\) in the allowable regime and for some finite constant \(C(\Lm)\).
This can be implemented through spectral norm regularization, contractive constraints, or dissipative terms in the physical realization of the PITN.

Semantically, \(\Lm\) ensures that the growth of complexity is tempered: the system can branch and explore without degenerating into divergence or collapse.

\section{Fractal Topology of Knowledge}

\subsection{The prime-word tree}

The set of all finite prime words
\[
\mathscr{W} = \bigcup_{t \geq 0} \mathbb{P}^t
\]
forms a rooted, infinitely branching tree.
The empty word \(\epsilon\) corresponds to the identity operator \(\Acal_{\epsilon} = \mathrm{Id}\).
Each word \(w\) has children \((w,p)\) for each prime \(p\), and the tree has the natural prefix order.

The mapping
\[
w \mapsto \Acal_{w}
\]
embeds this combinatorial tree into the operator space, and via the action of \(\Acal_w\) on \(\Mcal\), into the state space itself.
Different branches generally correspond to different semantic and dynamical regimes.

\subsection{Self-similarity and scale invariance}

Because the same finite alphabet (subsets of \(\mathbb{P}\)) is reused at every depth, subtrees rooted at different words are structurally similar.
Under suitable parameter-sharing or architectural constraints on the operators \(A_p\), this leads to a form of self-similarity: local dynamics around a word \(w\) can resemble the global dynamics on a rescaled manifold.

This fractal structure supports:
\begin{itemize}[nosep]
    \item reuse of learned local motifs (prime subsequences) across domains,
    \item hierarchical composition of skills by extending words,
    \item recursive embedding of tasks at different levels of abstraction.
\end{itemize}

\subsection{Dimensional unfolding via \texorpdfstring{\(\beta(t)\)}{beta(t)}}

Introduce a meta-learning variable \(\beta(t)\), which controls the set of primes that are \emph{active} at time \(t\).
For instance, define
\[
\mathbb{P}_{\beta(t)} := \{p \in \mathbb{P} : f(p) \leq \beta(t)\},
\]
for some monotonically increasing function \(f\) (e.g., \(f(p) = \log p\)).
At early times (small \(\beta\)), only small primes are available; as \(\beta\) grows, larger primes---and thus more complex operators---are unlocked.

The effective dimensionality of the reachable region in \(\Mcal\) grows with \(\beta\), because new operators open new directions in the manifold's tangent space.
However, previously established structures, anchored by \(\Lm\) and the existing words, remain intact.
This realizes a form of \emph{dimensional unfolding}: the knowledge manifold expands while preserving foundational structure.

\section{Hierarchical Decision-Making on the Fractal Tree}

\subsection{Policy over prime words}

Consider an agent interacting with an environment.
At each time step \(t\), the agent observes \(o_t\), which may depend on the environment state and the internal configuration \(x_t\).
The agent maintains a prime word \(w(t)\), and thus an operator \(\Awt\) and state \(x_t = \Awt(x_0)\).

A policy \(\pi\) in NAFL is a conditional distribution over primes:
\[
\pi(p \mid o_t, w(t)) = \Prob(p_{i_{t+1}} = p \mid o_t, w(t)).
\]
Upon sampling \(p_{i_{t+1}}\), the word extends to
\[
w(t+1) = (w(t), p_{i_{t+1}}),
\]
the composite operator updates to
\[
\Awp = A_{p_{i_{t+1}}} \Awt,
\]
and the configuration becomes
\[
x_{t+1} = \Awp(x_0).
\]

\subsection{Value functions over composite operators}

Define a reward signal \(r_t\) and a discount factor \(\gamma \in (0,1)\).
The value of a word \(w(t)\) under policy \(\pi\) is
\[
V^\pi(w(t)) = \E_\pi\left[ \sum_{k=t}^{\infty} \gamma^{k-t} r_k \,\middle|\, \Awt \right].
\]
Similarly, define an action-value function over candidate primes:
\[
Q^\pi(w(t), p) = \E_\pi\left[ \sum_{k=t}^{\infty} \gamma^{k-t} r_k \,\middle|\, w(t+1) = (w(t), p) \right].
\]

Decision-making is then the problem of choosing primes that extend the word so that the resulting composite operators land in high-value, stable regions of \(\Mcal\).

\subsection{Sub-task factorization as word decomposition}

A complex macro-goal can be represented as a constraint on acceptable words or operators, e.g. a set \(\mathscr{G} \subset \mathscr{W}\) of goal words whose operators map the current configuration into target regions of \(\Mcal\).

Sub-tasks emerge as:
\begin{itemize}[nosep]
    \item prefixes of goal words,
    \item recurrent motifs (subwords) with robust positive contributions to value,
    \item local branches that realize particular semantic transformations.
\end{itemize}

Hierarchical plans are then sequences of prime subsequences, assembled by navigating the prime-word tree.

\subsection{Contextual path traversal and ``fractal rollouts''}

To evaluate potential decisions, the agent can simulate candidate extensions \(\Delta w = (p_{i_{t+1}}, \dots, p_{i_{t+\tau}})\) and approximate the resulting composite operators \(\Acal_{(w(t), \Delta w)}\).
Because the tree is self-similar, local models learned in one sub-tree can be reused in structurally similar subtrees, enabling efficient generalization.

This process of simulating paths down the tree---rather than simply simulating in a flat latent space---constitutes \emph{fractal rollouts}.
They provide a natural mechanism for multi-step lookahead that respects the compositional structure of knowledge.

\subsection{Ecological entanglement via \texorpdfstring{\(Q_{\mathrm{ent}}\)}{Q\_ent}}

Let \(Q_{\mathrm{ent}}\) be an entanglement functional defined on pairs (or sets) of words:
\[
Q_{\mathrm{ent}} : \mathscr{W} \times \mathscr{W} \to \R.
\]
Intuitively, \(Q_{\mathrm{ent}}(w_a, w_b)\) measures how strongly the branches rooted at \(w_a\) and \(w_b\) are coupled in terms of semantics, dynamics, or global constraints.

In ecological decision-making, a change along one branch (e.g., resource allocation) must respect constraints on another branch (e.g., security, fairness, or long-term sustainability).
This can be modeled by including a coupling term in the objective:
\[
\Lcal_{\mathrm{total}} = \Lcal_{\mathrm{RL}} + \lambda_{\mathrm{ent}} \sum_{(w_a,w_b)} Q_{\mathrm{ent}}(w_a, w_b),
\]
where \(\Lcal_{\mathrm{RL}}\) is a standard RL loss (e.g., policy gradient or actor--critic), and \(\lambda_{\mathrm{ent}}\) regulates the strength of cross-branch coupling.

The functional \(Q_{\mathrm{ent}}\) can depend on:
\begin{itemize}[nosep]
    \item structural similarity (shared subsequences),
    \item shared operators or parameterizations,
    \item empirical co-variations in rewards or constraints.
\end{itemize}

\section{Advantages over Classical RL and Hierarchical RL}

\subsection{State representation}

Classical RL usually assumes a Markovian state \(s_t\) that contains all relevant information for decision-making.
While history can be encoded into \(s_t\) via recurrent networks, the underlying operator algebra is typically treated as associative and the history is compressed.

In contrast, NAFL defines the state as the pair \((x_t, w(t))\), with the operator \(\Awt\) being central.
Different histories correspond to different operators even if \(x_t\) coincides, and this difference matters for future evolutions.

\subsection{Hierarchy building}

Hierarchical RL introduces options, skills, or sub-goals, but often requires manual decomposition or heuristic discovery.
In NAFL, hierarchy is a native structural property of the prime-word tree:
\begin{itemize}[nosep]
    \item every prefix is a higher-level context,
    \item every extension is a more specialized sub-skill,
    \item recurring subwords become emergent options without additional machinery.
\end{itemize}

\subsection{Catastrophic forgetting}

Adding new tasks in classical deep RL often overwrites existing weights, leading to catastrophic forgetting.
Mitigation strategies (e.g., replay buffers, regularization) only partially address the issue.

In NAFL, new knowledge corresponds to:
\begin{itemize}[nosep]
    \item exploring new branches of the prime-word tree, or
    \item refining operators at specific words.
\end{itemize}
Because branches are structurally separated and anchored by \(\Lm\), the risk of overwriting foundational paths is reduced.
Stability can be further enforced by regularizing parameter updates on well-established words and confining large updates to less-constrained branches.

\section{Minimal Implementation Sketch}

This section provides a lightweight pseudocode sketch of a ``toy NAFL'' system, intended as a blueprint rather than an optimized implementation.

\subsection{Core objects}

\subsubsection{Prime operator library}

Each prime \(p\) has an associated parameterized operator \(A_p\).
In practice, one could map primes to indices of a small operator library.

\begin{lstlisting}[language=Python, caption={Toy prime operator library interface}]
class PrimeOperator:
    def __init__(self, prime_id, params):
        self.prime = prime_id
        self.params = params  # e.g., weights of a small neural block

    def __call__(self, x):
        # Apply the operator to state x
        # e.g., x' = f(x; params)
        return apply_tensor_block(x, self.params)

class PrimeOperatorLibrary:
    def __init__(self, prime_list):
        self.ops = {p: PrimeOperator(p, init_params_for_prime(p))
                    for p in prime_list}

    def apply(self, prime, x, Lambda_m):
        op = self.ops[prime]
        x_next = op(x)
        return Lambda_m * x_next
\end{lstlisting}

\subsubsection{Prime word and trie}

The prime word is stored explicitly as a list of primes, and a trie is used for efficient lookup of prefixes and subwords.

\begin{lstlisting}[language=Python, caption={Prime word and trie structures}]
class PrimeWord:
    def __init__(self, primes=None):
        self.primes = list(primes) if primes is not None else []

    def extend(self, p):
        new_word = PrimeWord(self.primes + [p])
        return new_word

class WordTrieNode:
    def __init__(self):
        self.children = {}
        self.metadata = {}  # e.g., value estimates, visit counts

class WordTrie:
    def __init__(self):
        self.root = WordTrieNode()

    def insert(self, word, metadata_update=None):
        node = self.root
        for p in word.primes:
            if p not in node.children:
                node.children[p] = WordTrieNode()
            node = node.children[p]
        if metadata_update is not None:
            node.metadata.update(metadata_update)

    def get_node(self, word):
        node = self.root
        for p in word.primes:
            if p not in node.children:
                return None
            node = node.children[p]
        return node
\end{lstlisting}

\subsubsection{Path-conditioned policy}

The policy network conditions on both the observation and an encoding of the current prime word.

\begin{lstlisting}[language=Python, caption={Path-conditioned policy skeleton}]
class NAFLPolicy:
    def __init__(self, prime_list):
        self.prime_list = prime_list
        self.policy_params = init_policy_params()

    def encode_word(self, word):
        # e.g., simple embedding of primes followed by a small RNN
        return encode_prime_sequence(word.primes)

    def __call__(self, observation, word):
        word_encoding = self.encode_word(word)
        logits = policy_network_forward(observation, word_encoding,
                                        self.policy_params)
        return logits  # to be turned into a distribution over primes

    def sample_prime(self, observation, word):
        logits = self(observation, word)
        return sample_from_logits_over_primes(logits, self.prime_list)
\end{lstlisting}

\subsection{Learning loop outline}

A simplified learning loop might look like:

\begin{lstlisting}[language=Python, caption={Simplified NAFL training loop sketch}]
Lambda_m = init_Lambda_m()
prime_list = [2, 3, 5, 7, 11]  # small set for toy experiments

library = PrimeOperatorLibrary(prime_list)
trie = WordTrie()
policy = NAFLPolicy(prime_list)

for episode in range(num_episodes):
    x = init_state()
    word = PrimeWord()
    observation = env.reset()

    for t in range(max_steps):
        # Choose next prime
        p_next = policy.sample_prime(observation, word)

        # Extend word and apply operator
        word_next = word.extend(p_next)
        x = library.apply(p_next, x, Lambda_m)

        # Step environment (e.g., based on derived action from x)
        action = derive_action_from_state(x)
        observation, reward, done, info = env.step(action)

        # Store transition (word, p_next, reward, word_next) in memory
        memory.store(word, p_next, reward, word_next)

        word = word_next
        if done:
            break

    # After episode: update policy and operators using stored transitions
    update_policy_and_operators(memory, policy, library, trie, Lambda_m)
\end{lstlisting}

This sketch omits many important details (value estimation, entropy regularization, stability constraints), but it shows how the key NAFL concepts---prime operators, prime words, and path-conditioned policies---can be instantiated in a concrete learning system.

\section{Discussion and Future Work}

Non-Associative Fractal Learning reframes cognition as navigation in a non-associative, prime-indexed operator space.
The path---the ordered sequence of prime moves---is the memory, and the induced fractal topology supports hierarchical, context-sensitive decision-making.

Several directions for further development include:
\begin{itemize}[nosep]
    \item introducing explicit non-associative bracketing degrees of freedom,
    \item exploring different realizations of \(\Mcal\) (e.g., quantum tensor networks, differentiable manifolds),
    \item investigating learning dynamics under various choices of \(\Lm\) and \(\beta(t)\),
    \item designing benchmark environments where NAFL's path-dependent advantages can be empirically quantified,
    \item formalizing \(Q_{\mathrm{ent}}\) using tools from information theory, category theory, or operator algebras.
\end{itemize}

As the mathematical and empirical theory matures, NAFL may provide a principled, scalable way to integrate multiplicity-based models of physics, computation, and cognition into a single AGI-capable architecture.

\nocite{*}
\bibliographystyle{unsrt}  
\bibliography{references}
\end{document}